\section{The pressure-consistent method}\label{sec:consistent}

\begin{proposition}[The naive consistent form, and its $L^2_0$ variant]\label{prop:singular}
\begin{enumerate}[label=(\alph*)]
\item Replace $\ip{p_h-\pbG(p_h)}{v\cdot n}$ by $\ip{p_h}{v\cdot n}$. The resulting linear system is square ($\VR\times\Pih$ tested by $\VR\times\Pih$). It has the kernel vector $(0,1)$ and is therefore singular. Its left null space is therefore nontrivial as well, and the system is solvable only for data satisfying one linear compatibility condition.
\item \emph{The $L^2_0$ variant.} Restrict the pressure to $\Pih\cap L^2_0(\Oh)$ and test the continuity row only with $q\in\Pih\cap L^2_0(\Oh)$, keeping the momentum row $N_h(u_h,v)-(p_h,\div v)+\ip{p_h}{v\cdot n}=(\ft,v)$ for all $v\in\VR$. Any solution has $\div u_h\equiv c$, a constant, with
\[
c\,|\Oh|=\int_{\partial R}\gI\cdot\nu+\int_{\Gh}u_h\cdot n ,
\]
$n$ being the outward normal of $\Oh$. If $c=0$, then $(u_h,p_h)$ solves the naive system of (a).
\item Let $\gamma\ge\gamma_2$ (Theorem~\ref{thm:D}). The naive system of (a) is solvable if and only if the $\CNS$ pressure satisfies $\pbG(p_h^{\CNS})=0$, and then its solutions are $(u_h^{\CNS},p_h^{\CNS}+\text{const})$. If moreover the $L^2_0$ variant is uniquely solvable, then $c=0$ if and only if the naive system is solvable.
\end{enumerate}
\end{proposition}

\begin{proof}
(a) For $v\in\VR$,
\[
-(1,\div v)+\ip{1}{v\cdot n}=-\int_{\partial\Oh}v\cdot n+\int_{\Gh}v\cdot n=0 ,
\]
and the continuity row vanishes at $u=0$. A square matrix with a nontrivial kernel has a nontrivial left kernel.

(b) The momentum rows of (a) and (b) coincide. Since $\div u_h\in\Pih$ (Lemma~\ref{lem:divrange}) and $(q,\div u_h)=0$ for all mean-free $q\in\Pih$, $\div u_h$ is constant; integrate over $\Oh$, using $u_h=\gI$ on $\partial R$. If $c=0$, then $(q,\div u_h)=0$ for all $q\in\Pih$.

(c) Write $B^\ast$ as in the proof of Theorem~\ref{thm:D}. For every $r\in\Pih$ the naive pressure form satisfies $-(r,\div v)+\ip{r}{v\cdot n}=B^\ast(r-\pbG(r),v)$, and it is invariant under $r\mapsto r+\text{const}$. If $(u,p)$ solves the naive system, then $(u,p-\pbG(p))$ solves $\CNS$, so by uniqueness (Theorem~\ref{thm:D}) $u=u_h^{\CNS}$ and $p-\pbG(p)=p_h^{\CNS}$, which has zero mean on $\Gh$. Conversely, if $\pbG(p_h^{\CNS})=0$, the two pressure forms agree at $p_h^{\CNS}$. For the last statement: if the naive system is solvable, shift its pressure into $L^2_0$; this gives a solution of the $L^2_0$ variant with $\div u_h=0$, which by uniqueness is the solution, so $c=0$. The converse is (b).
\end{proof}

The singularity is the analogue, for our square Scott--Vogelius system, of the ill-posedness that Liu, Neilan and Otus note for their boundary-multiplier method when the zero-flux constraint is omitted \cite[Remark~3.2]{LiuNeilanOtus2023}. The $L^2_0$ variant is the standard way to fix the constant; in the fully weak setting of Frachon, Nilsson and Zahedi the same computation shows that $\div u_h$ equals the multiplier of the mean constraint \cite[(5.7)--(5.9)]{FrachonNilssonZahedi2024}. With weak imposition nothing forces $\int_{\Gh}u_h\cdot n=-\int_{\partial R}\gI\cdot\nu$, even for flux-corrected data, so generically $c\ne0$ and the $L^2_0$ variant is \emph{not} exactly divergence-free. Unique solvability of the $L^2_0$ variant reduces to the nonvanishing of a single scalar (for zero data, write $u_h=u_0+c\,w_c$ with $\div w_c\equiv1$; then $(u_0,p_h-\pbG(p_h))$ is $c$ times a fixed $\CNS$ solution, and the constraint $\pbG(\cdot)=0$ on its pressure must force $c=0$). It is generic and held in every computation below, but we do not prove it.

Numerically ($k=4$, meshes of Section~\ref{sec:numerics}), the kernel of (a) is exactly one-dimensional: in a dense SVD of the naive system for test B with $\mu=20$ and an element-wise $L^2$-orthonormal pressure basis, at $N=8$ and $12$ the smallest singular values are $8\cdot10^{-17}$ and $3\cdot10^{-16}$, the next ones $2.4\cdot10^{-2}$ and $8.1\cdot10^{-3}$. For the $L^2_0$ variant with $\mu=100$ (SuperLU, relative residual about $6\cdot10^{-16}$), $\div u_h$ is constant to round-off. For test A at $N=16$ it has $c=9.69\cdot10^{-4}$, so $\norm{\div u_h}_{L^2(\Oh)}=|c|\,|\Oh|^{1/2}=4.54\cdot10^{-3}$, and the $\CNS$ pressure has $\pbG(p_h)=7.5\cdot10^{-2}\ne0$. For test B at $N=16$ and $32$, $c=-5.98\cdot10^{-5}$ and $-2.94\cdot10^{-5}$; there the $L^2_0$ variant and $\CNS$ have $H^1$ velocity errors that agree to four digits.

Earlier versions of this paper reported that an iterated-penalty solver and a bordered direct solve of the naive form returned discrete divergence ``about $4.5\cdot10^{-3}$'', and read this as a sign of incompatibility. A bordered solve that fixes the pressure mean computes exactly the $L^2_0$ variant of (b), and the value coincides with $|c|\,|\Oh|^{1/2}=4.54\cdot10^{-3}$ for test A, $N=16$, $\mu=100$. The reported divergence is therefore the constant $c$ of (b), not a symptom of an incompatible system. We have not re-run the original solves, so the identification of the test case rests on this numerical coincidence.
%% VERIFY: the original 4.5e-3 runs (test, N, mu) were not archived; identification with test A, N=16, mu=100 is by value only (notes/v5/lit_report.md).

Let $\beta$ be the uniform inf-sup constant of Lemma~\ref{lem:infsup}.

\begin{theorem}[$\CNS$: well posed, with Scott's rate]\label{thm:D}
There is $\gamma_2\asymp\max(\gamma_0,\beta^{-2})$ such that for $\gamma\ge\gamma_2$ the method $\CNS$ has a unique solution and
\[
\tnorm{\ut-u_h}\le C\Bigl[(1+\gamma^{1/2})\hG^{3/2}+\eps\Bigr].
\]
Consequently, for every fixed $k\ge4$, for $\gamma$ in a bounded range $[\gamma_2,\gamma_{\max}]$ and $\hG\ge\hO^{2(\sigma_k-k)}$ (that is, $\hG\ge\hO^4$ for even $k$ and $\hG\ge\hO^2$ for odd $k$),
\[
\norm{\ut-u_h}_{H^1(\Oh)}\le C\bigl(\hG^{3/2}+\hO^k\bigr).
\]
This condition ensures that the flux term of $\eps$ is at most $Ch^k$. Without it, that single extra term is kept, and by the remark after Lemma~\ref{lem:approx} it is then genuinely present for the plain interpolant $\gI$. With the flux-corrected data $\tilde g_I$ of Definition~\ref{rs:def} no proviso is needed: $\norm{\ut-u_h}_{H^1(\Oh)}\le C(\hG^{3/2}+\hO^k)$ for every $\gamma\ge\gamma_2$, with $C$ independent of $\gamma$ and no relation between $\hG$ and $\hO$ (Theorem~\ref{rs:Dfc}). This is the rate Scott conjectured, for general Stokes data. The $L^2$ part of the $H^1$ norm uses the lift $E$ from the proof of Theorem~\ref{thm:B}(ii).
\end{theorem}

\begin{proof}
\begin{enumerate}
\item \emph{Consistency.} Write $B^\ast(q,v):=-(q,\div v)+\ip{q-\pbG(q)}{v\cdot n}$.
\begin{itemize}
\item For a constant $c$, $B^\ast(q+c,v)=B^\ast(q,v)-c\int_{\Gh}v\cdot n$.
\item $B^\ast(\pt,v)$ equals the pressure form of Lemma~\ref{lem:erroreq} minus $\pbG(\pt)\int_{\Gh}v\cdot n$.
\item Hence $p^\ast:=\pt-\pbG(\pt)$ satisfies $N_h(\ut,v)+B^\ast(p^\ast,v)=(\ft,v)+\Gc(v)$ for all $v\in\VR$.
\end{itemize}
With $e_u:=\ut-u_h$ and $e_p:=p^\ast-p_h$ this gives
\begin{equation}\label{eq:cnserr}
N_h(e_u,v)+B^\ast(e_p,v)=\Gc(v)\qquad\forall v\in\VR .
\end{equation}
\item \emph{Split.} Write $e_p=\eta+\xi$ with $\eta:=p^\ast-\pi_hp^\ast$ and $\xi\in\Pih$. For $v\in\VO$, $(\eta,\div v)=0$ because $\div\VO\subset\Pih$.
\item \emph{Pressure.} For $v\in\VO$, \eqref{eq:cnserr} gives
\[
(\xi-\bar\xi,\div v)=N_h(e_u,v)-\Gc(v)=a_h(e_u,v)-\ip{e_u}{\dn v}-\Gc(v).
\]
Using $(\rho/\mu)^{1/2}\le C$ and Lemma~\ref{lem:Gbound}, the right-hand side is at most $C(\tnorm{e_u}+\hG^{3/2})\norm{\nabla v}$. Lemma~\ref{lem:infsup} then gives $\beta\norm{\xi-\bar\xi}\le C(\tnorm{e_u}+\hG^{3/2})$.
\item \emph{Velocity.} For $z\in\Wh$, $e_h:=z-u_h\in\Zh$. Test \eqref{eq:cnserr} with $v=e_h$. Since $(e_p,\div e_h)=0$ and $\ip{\pbG(e_p)}{e_h\cdot n}=0$,
\[
c_1\tnorm{e_h}^2\le ME_h\tnorm{e_h}+|\Gc(e_h)|+|\ip{\eta}{e_h\cdot n}|+|\ip{\xi-\bar\xi}{e_h\cdot n}|,
\]
where
\begin{align*}
|\ip{\eta}{e_h\cdot n}|&\le C\hG^k(h/\mu)^{1/2}\tnorm{e_h},\\
|\ip{\xi-\bar\xi}{e_h\cdot n}|&\le C_I(\rho/h)^{1/2}\norm{\xi-\bar\xi}\,(h/\mu)^{1/2}\tnorm{e_h}\le C(c_0\gamma)^{-1/2}\norm{\xi-\bar\xi}\,\tnorm{e_h}.
\end{align*}
\item \emph{Absorb.} Insert step~3 and $\tnorm{e_u}\le E_h+\tnorm{e_h}$. If $\gamma\ge\gamma_2$ with $C(c_0\gamma_2)^{-1/2}\beta^{-1}\le c_1/2$, the pressure term is absorbed.
\item \emph{Uniqueness.} Apply steps 3--5 to the homogeneous problem ($\gI=0$, so $z=0\in\Wh$). This gives $u_h=0$, and then $p_h=c$ is constant. Now \eqref{eq:cnserr} reads $-c\int_{\Gh}v\cdot n=0$ for all $v$, and the edge bubble of Lemma~\ref{lem:divrange} gives $c=0$. The system is square, so uniqueness gives existence.\qedhere
\end{enumerate}
\end{proof}

\begin{remark}[Relation to the zero-flux formulation of Frachon, Nilsson and Zahedi]\label{rem:fnz}
Let $V_h^{R,\flat}:=\{v\in\VR:\int_{\Gh}v\cdot n=0\}$, which has codimension one (edge bubble, Lemma~\ref{lem:divrange}). Transcribed to our setting, the device of \cite[Section~5.1]{FrachonNilssonZahedi2024} tests the momentum row with the naive form $-(p_h,\div v)+\ip{p_h}{v\cdot n}$ only for $v\in V_h^{R,\flat}$, keeps the full continuity row, and fixes the pressure constant separately. On $V_h^{R,\flat}$ the naive and the $\CNS$ pressure forms coincide, because $\ip{\pbG(p_h)}{v\cdot n}=\pbG(p_h)\int_{\Gh}v\cdot n=0$; and $B^\ast(p+c,w)=B^\ast(p,w)-c\int_{\Gh}w\cdot n$ for the remaining direction $w$. Hence, for $\gamma\ge\gamma_2$: the $\CNS$ solution solves the zero-flux formulation; conversely, if $(u,p)$ solves the homogeneous zero-flux formulation, then $(u,p+c)$ solves homogeneous $\CNS$ for the $c$ that satisfies the $w$-row, so $u=0$. So the two formulations have the \emph{same velocity}, and pressures that differ by a constant. $\CNS$ is a square packaging of the same device whose constant is fixed automatically, with $p_h\approx\pt-\pbG(\pt)$.
\end{remark}

\subsection{A matching lower bound: the rate \texorpdfstring{$\hG^{3/2}$}{h^{3/2}} is sharp}\label{lb:sec}

Theorem~\ref{thm:D} gives $\norm{\nabla(\ut-u_h)}\le\tnorm{\ut-u_h}\le C[(1+\gamma^{1/2})\hG^{3/2}+\eps]$ for $\CNS$, and Theorem~\ref{thm:A} gives the same bound for $\GS$ when $G=0$. Here we prove the reverse inequality in the $H^1$ seminorm whenever the wall shear is not identically zero.

\emph{Status of the hypothesis.} The proof uses a local condition (M3) on the boundary stars (Assumption~\ref{lb:ass}). For $k\ge7$ a single-triangle field gives it on every mesh satisfying (M0)--(M2) (Lemma~\ref{lb:k7}). For every $k\ge4$ it is proved in Section~\ref{m3:sec} (Theorem~\ref{m3:thm}) under (M0)--(M2) and the smallness condition (S): $\hG\le2\sin(\theta_{\min}/4)$, with an explicit $P_4$ witness on three triangles of each boundary star, a positive constant $\kappa_0$ by compactness, and certified numerical values for two angle bounds. So the lower bound is unconditional for every $k\ge4$. (Earlier versions of this paper proved it only for $k\ge7$ and verified (M3) numerically for $k=4$.) The proof does \emph{not} go through the transposed-gradient term of Lemma~\ref{lem:transposed}; see Section~\ref{lb:sec:route} for why that route fails for $\CNS$.

\subsubsection{Notation}
On $\Gamma$ let $n$ be the unit normal pointing into $D$ and $\tau$ the unit tangent, oriented so that $\tau(x^\ast)\cdot\tau_e>0$ for $x\in e$ ($\tau_e$ is $n_e$ rotated in the same sense). By Lemma~\ref{lem:wall}, $\dn u=W\tau$ on $\Gamma$, where
\[
W:=\dn u\cdot\tau\qquad\text{(the wall shear)}.
\]
For $e\in\EG$, $T_e$ is the triangle with edge $e$, $s$ is arclength on $e$ from its midpoint $m_e$, and $W_e:=W(m_e^\ast)$. No triangle has two edges on $\Gh$: two such edges would share a vertex $z$ of $\Gh$, and the angle of $\Oh$ at $z$ is $\pi+O(\hG)>\pi$. Define
\[
Y_h:=\Zh\cap\VO,\qquad Y_h^1:=\Bigl\{v\in Y_h:\ \int_e\dn v=0\ \ \forall e\in\EG\Bigr\}.
\]

\subsubsection{The local assumption}
Write $\mathcal M(v):=-\tfrac12\sum_{e\in\EG}\int_e s^2\,\dn v\cdot\tau_e\,ds$. For a vertex $z$ of $\Gh$, let $e_z$ be the edge of $\Gh$ that follows $z$ counterclockwise; $z\mapsto e_z$ is a bijection between the vertices and the edges of $\Gh$.

\begin{assumption}[Boundary test fields]\label{lb:ass}
\textbf{(M3)} There are constants $\kappa_0>0$ and $K_{\mathrm{ov}},C_s$ such that every vertex $z$ of $\Gh$ carries a field $v_z\in Y_h^1$ with the following properties:
\begin{enumerate}[label=(\roman*)]
\item $\norm{\nabla v_z}=1$;
\item $v_z$ is supported in a patch $\omega_z$ of triangles within distance $C_s\hG$ of $z$, and every triangle lies in at most $K_{\mathrm{ov}}$ patches;
\item $\mathcal M(v_z)\ge\kappa_0\,|e_z|^2$, with $e_z$ as above.
\end{enumerate}
\end{assumption}

Condition (iii) is invariant under scaling, since $\norm{\nabla v}$ is scale invariant in two dimensions. Only the sign of $v_z$ is chosen in (iii): $\mathcal M(-v)=-\mathcal M(v)$.

\begin{lemma}[(M3) for $k\ge7$]\label{lb:k7}
If $k\ge7$, (M3) holds with $\omega_z=T_{e_z}$. Here $K_{\mathrm{ov}}=1$ and $\kappa_0$ depends only on (M0).
\end{lemma}

\begin{proof}
Let $T=T_{e_z}$, and let $\lambda_1,\lambda_2,\lambda_3$ be its barycentric coordinates with $\lambda_3=0$ on $e:=e_z$. Put $b:=\lambda_1\lambda_2\lambda_3$ and $\sigma:=2s/|e|=\lambda_1-\lambda_2$ on $e$ (up to sign), and set
\[
\phi:=b^2\bigl((\lambda_1-\lambda_2)^2-\tfrac17\bigr)\in P_8,\qquad v:=\curl\phi\in P_7^2 .
\]
\begin{itemize}
\item $\phi$ vanishes to second order on $\partial T$, so $v=0$ on $\partial T$. Extended by zero, $v$ is continuous, lies in $\VO\subset\VR$ (because $k\ge7$), and satisfies $\div v=0$. Hence $v\in Y_h$.
\item On $e$, $\dn v\cdot\tau_e=\pm\partial_{nn}\phi=\pm2|\nabla\lambda_3|^2(\lambda_1\lambda_2)^2(\sigma^2-\tfrac17)=\pm\tfrac18|\nabla\lambda_3|^2(1-\sigma^2)^2(\sigma^2-\tfrac17)$.
\item Since $\int_{-1}^1(1-\sigma^2)^2=\frac{16}{15}$ and $\int_{-1}^1\sigma^2(1-\sigma^2)^2=\frac{16}{105}$, this has mean zero. With $\dn v\cdot n_e=0$ (Lemma~\ref{lb:normal} below), $\int_e\dn v=0$, so $v\in Y_h^1$.
\item Its second moment is $\int_{-1}^1\sigma^2(1-\sigma^2)^2(\sigma^2-\frac17)\,d\sigma=\frac{16}{315}-\frac{16}{735}=\frac{64}{2205}\neq0$.
\end{itemize}
The other two edges of $T$ are interior edges and do not enter $\mathcal M$. Hence $\mathcal M(v)\ne0$. The ratio $|\mathcal M(v)|/(|e|^2\norm{\nabla v})$ is invariant under similarity. It is a continuous, nonvanishing function of the shape of $T$, and the shapes allowed by (M0) form a compact set, so the ratio is bounded below by some $\kappa_0>0$. Normalise and choose the sign.
\end{proof}

\begin{remark}[$k=4$]\label{lb:k4}
For $k=4$ no field in $Y_h^1$ can be supported in a single triangle, or in the two triangles of one first-layer quadrilateral; for the production meshes the latter space is $\{0\}$ (Section~\ref{lb:sec:num}). The smallest natural patch is the boundary star $\omega_z$ (the triangles at $z$).

Write $v=\curl\phi$ with $\phi$ a $C^1$ piecewise quintic; since $v=0$ on $\partial\omega_z$ and on $\Gh$, we may take $\phi=0$ and $\nabla\phi=0$ on $\partial\omega_z$ and on the two edges $e,e'$ of $\Gh$ at $z$. On such an edge $e$, $\dn v\cdot\tau_e=\pm\partial_{n_en_e}\phi|_{T_e}$ is a cubic with zero mean ($v\in Y_h^1$). At the far endpoint of $e$ the gradient of $\phi|_{T_e}$ vanishes along two independent lines, so its Hessian vanishes there. If also $\partial_{n_en_e}\phi|_{T_e}(z)=0$, the cubic vanishes at both ends of $e$ and has zero mean, so it is odd about the midpoint and $\int_es^2\,\dn v\cdot\tau_e=0$. Hence (M3) requires $\partial_{n_en_e}\phi|_{T_e}(z)\ne0$ for at least one of the two edges $e$ at $z$, the value being taken from the triangle $T_e$. This is possible only because $C^1$ quintics can have Hessians that jump across the interior edges at $z$: an Argyris ($C^2$-vertex) stream function has a single Hessian at $z$, which vanishes on both boundary directions and hence is zero. Such jumps exist at a boundary vertex with at least three triangles, which is (M1$'$).

Section~\ref{m3:sec} proves (M3) for $k=4$ (hence every $k\ge4$) by constructing such a field explicitly. As a check, the floating-point star computation of the earlier versions gives:
\begin{itemize}
\item for every boundary vertex of every production mesh ($N=8$, $12$, $16$, $20$, $24$, $32$, $48$, $64$, $96$, $128$, $192$, $256$) and also for $N=512$ and $1024$ (Table~\ref{lb:tab:stars}): $\dim Y^1(\omega_z)=1$, $\min_z\kappa_z\ge0.00288$ in every case, and $\min_z\kappa_z$ converges to about $0.0029$ as $N$ grows;
\item for randomly generated three- and four-triangle boundary stars: $\dim Y^1(\omega_z)=2m-5$ for $m$ triangles, and $\kappa_z>0$ in every sample.
\end{itemize}
(M3) is monotone in the degree: a field in $Y_h^1$ for degree $k$ is continuous piecewise $P_{k'}$ for every $k'\ge k$, and $\mathcal M$ and $\norm{\nabla\cdot}$ do not depend on $k$. So (M3) for $k=4$ on a mesh implies (M3) for every $k\ge4$ on the same mesh. The exact star constants of Section~\ref{m3:sec:mesh} reproduce the table below to all printed digits.
\end{remark}

% (M3) proof, from notes/v4/m3_proof.tex as audited in notes/v6/m3 and notes/v6/referee3_ns_m3.md.
\subsubsection{Proof of (M3) for every \texorpdfstring{$k\ge4$}{k>=4}}\label{m3:sec}

This subsection removes the hypothesis (M3) of Theorem~\ref{lb:thm} for $4\le k\le 6$. The construction was first drafted in version~4 of this work; the two geometric steps marked \emph{exact} below replace ``for $\hG$ small'' arguments of that draft, and the witness identities were re-verified independently (\texttt{notes/v6/m3/verify\_witness.py}). Throughout, $z$ is a vertex of $\Gh$, and the triangles at $z$ form a fan $T_j=(z,p_{j-1},p_j)$, $j=1,\dots,m$, listed counterclockwise, with $[z,p_0]$ and $[z,p_m]$ the two edges of $\Gh$ at $z$. We write $\alpha_j$ for the angle of $T_j$ at $z$, $\beta_j$ and $\gamma_j$ for its angles at $p_{j-1}$ and $p_j$, and $\Theta_z=\sum_j\alpha_j$ for the angle of $\Oh$ at $z$. Put $z=0$, and for points $a,b$ let $a\times b:=a_1b_2-a_2b_1$ and $D_{ij}:=p_i\times p_j$.

\paragraph{Hypotheses.} Besides (M1$'$) we use:
\begin{itemize}
\item[(M0$_\theta$)] every triangle with a vertex on $\Gh$ has all its angles $\ge\theta_{\min}$, where $0<\theta_{\min}\le\pi/3$;
\item[(S)] $\hG\le2\sin(\theta_{\min}/4)$.
\end{itemize}
(M0$_\theta$) is the quantitative form of the shape regularity in (M0), and (S) is a smallness condition of the type allowed by (D). (M2) is used only through (M1$'$).

\begin{lemma}[Exact boundary geometry]\label{m3:geom}
\begin{enumerate}[label=(\alph*)]
\item At every vertex $z$ of $\Gh$, with $\ell_1,\ell_2$ the lengths of the two chords at $z$, $\Theta_z-\pi=\arcsin(\ell_1/2)+\arcsin(\ell_2/2)$. Under (S), $|\Theta_z-\pi|\le\theta_{\min}/2$, so the condition
\begin{itemize}
\item[(A$_\delta$)] $|\Theta_z-\pi|\le\delta$ at every vertex $z$ of $\Gh$
\end{itemize}
holds with $\delta=\theta_{\min}/2<\theta_{\min}$.
\item On every mesh, the only edges of $\Gh$ contained in the star of $z$ are $[z,p_0]$ and $[z,p_m]$:
\begin{itemize}
\item[(G)] no other edge of the star lies on $\Gh$.
\end{itemize}
\end{enumerate}
\end{lemma}

\begin{proof}
(a) A chord of length $\ell$ of the unit circle subtends the central angle $2\arcsin(\ell/2)$, and the angle between the chord and the tangent at either endpoint is half of it. The interior angle of $\Oh=R\setminus P_h$ at $z$ is $\pi$ plus the two tangent--chord angles. Under (S), $\arcsin(\ell_i/2)\le\arcsin(\sin(\theta_{\min}/4))=\theta_{\min}/4$.

(b) A triangle of the star with an edge on $\Gh$ not through $z$ has all three vertices on $\partial P_h$, so it lies in the convex polygon $\overline{P_h}$; its interior would then be disjoint from $\Oh$. An edge of $\Gh$ through $z$ is one of the two chords at $z$.
\end{proof}
\begin{theorem}[(M3) holds for $k\ge4$]\label{m3:thm}
Assume (M0$_\theta$), (M1$'$) and (A$_\delta$) with $0\le\delta<\theta_{\min}$, and let $k\ge4$. (By Lemma~\ref{m3:geom}, (A$_\delta$) holds with $\delta=\theta_{\min}/2$ under (S), and (G) holds on every mesh.) Then (M3) holds with $K_{\mathrm{ov}}=3$, $C_s=(\sin\theta_{\min})^{-3}$ and
\[
\kappa_0=\kappa_\ast(\theta_{\min},\delta)>0,
\]
where $\kappa_\ast$ is the minimum of the explicit continuous function $\kappa_w$ of Lemma~\ref{m3:norm} over the compact families $F_3(\theta_{\min},\delta)$ and $F_s(\theta_{\min},\delta)$ of Definition~\ref{m3:fam}. The patch $\omega_z$ is the star of $z$ if $m=3$, and the three triangles $T_1,T_2,T_3$ of the star next to $e_z$ if $m\ge4$. Moreover:
\begin{enumerate}[label=(\alph*)]
\item $\kappa_\ast(\theta_{\min},\delta)\ge\kappa_{\rm el}(\theta_{\min},\delta)$, the closed-form expression of Corollary~\ref{m3:el};
\item rigorous interval certificates give $\kappa_\ast(30^\circ,10^\circ)\ge0.0012$ and $\kappa_\ast(20^\circ,15^\circ)\ge0.0003$ (Section~\ref{m3:sec:cert}). The certificate parameters $(\theta_{\min},\delta)$ are separate from (S): the value at $(30^\circ,10^\circ)$ applies when $\theta_{\min}\ge30^\circ$ and $\sup_z|\Theta_z-\pi|\le10^\circ$, e.g.\ for $\hG\le2\sin5^\circ$;
\item on every mesh of Section~\ref{sec:numerics} ($N=8,\dots,1024$) every boundary star carries a field with $\kappa_z\ge0.00287$, for either choice of $e_z$ (Section~\ref{m3:sec:mesh}); these are the exact star constants of Table~\ref{lb:tab:stars}.
\end{enumerate}
\end{theorem}

Consequently Theorem~\ref{lb:thm} and Corollary~\ref{lb:cor} hold for every $k\ge4$ under (M0)--(M2), (D) and (S), with $\kappa_0=\kappa_\ast(\theta_{\min},\theta_{\min}/2)$. For $k\ge7$, Lemma~\ref{lb:k7} gives a single-triangle field as well. The witness below is a $P_4$ field, so by the monotonicity in Remark~\ref{lb:k4} it serves for all $k\ge4$.

\paragraph{The idea.} Remark~\ref{lb:k4} shows that a field in $Y_h^1$ can have $\mathcal M\ne0$ only if the Hessian of its stream function at $z$ jumps across the interior edges at $z$. On a fan of three triangles the homogeneous quadratic $C^1$ splines vanishing doubly on the two outer rays form a one-dimensional space, spanned by a spline whose Hessians are given by a classical Pl\"ucker-type identity; all its coefficients have a sign when $\alpha_1+\alpha_2<\pi$ and $\alpha_2+\alpha_3<\pi$. We complete this quadratic part by a cubic part, in closed form, to a $C^1$ quintic stream function with zero boundary data and zero means on the chords. Its second moment is a sum of two terms of the same sign.

\paragraph{Stream functions.}
Let $F$ be a fan of triangles at $z$ as above, a simply connected polygon. A continuous piecewise-$P_k$ field $v$ on $F$ with $v=0$ on $\partial F$ and $\div v=0$ is $v=\curl\phi$ with $\phi\in C^1(F)$ piecewise $P_{k+1}$ and $\phi=|\nabla\phi|=0$ on $\partial F$, and conversely; moreover $\norm{\nabla v}_F=\norm{D^2\phi}_F$. If $e=[z,p]\subset\partial F$, then $\partial_t\partial_n\phi=0$ on $e$, so $\dn v=(\partial_{nn}\phi)\,R n$ with $R$ the rotation by $-\pi/2$. Since $\tau_e$ is $n_e$ rotated in a fixed sense for all edges,
\begin{equation}\label{m3:dnv}
\dn v\cdot\tau_e=\sigma\,\partial_{nn}\phi\quad\text{on every edge }e\subset\partial F\cap\Gh,\qquad\sigma\in\{\pm1\}\text{ the same for all edges}.
\end{equation}
So $\int_e\dn v=0$ if and only if $\int_e\partial_{nn}\phi=0$ (the normal component vanishes by Lemma~\ref{lb:normal}), and $\mathcal M(v)=-\tfrac\sigma2\sum_e\int_e s^2\partial_{nn}\phi$.

\paragraph{$C^1$ quintics on a fan.}
On $T_j$ let $L_j$ be the barycentric coordinate of $z$; it is affine, $L_j(z)=1$, and $L_j=0$ on the line $p_{j-1}p_j$. Put $N_j(x):=p_j\times x$, a linear form vanishing on the ray through $p_j$.

\begin{lemma}[Gluing]\label{m3:glue}
Let $\phi$ be piecewise quintic on the fan with $\phi=|\nabla\phi|=0$ on its boundary, and $\phi_j:=\phi|_{T_j}$. Then $\phi_j=L_j^2(Q_j+C_j)$ with $Q_j$, $C_j$ homogeneous of degrees $2$ and $3$. On the boundary rays, $Q_1=aN_0^2$, $C_1=N_0^2\ell_1$, $Q_m=a'N_m^2$, $C_m=N_m^2\ell_m$ with $a,a'\in\R$ and $\ell_1,\ell_m$ linear. Write $L_{j+1}-L_j=\beta_jN_j$ ($\beta_j=0$ iff $p_{j-1},p_j,p_{j+1}$ are collinear). Then $\phi$ is $C^1$ across the ray through $p_j$ if and only if
\begin{enumerate}[label=(G\arabic*)]
\item $Q_j-Q_{j+1}\in\R N_j^2$;
\item $C_j(p_j)=C_{j+1}(p_j)$;
\item $\nabla(C_j-C_{j+1})(p_j)=2\beta_jQ_j(p_j)\,\nabla N_j$;
\item $\beta_j\bigl(Q_j(p_j)+C_j(p_j)\bigr)=0$.
\end{enumerate}
\end{lemma}

\begin{proof}
$\phi_j$ vanishes to first order on the line $p_{j-1}p_j$, so $L_j^2$ divides it; the quotient $q_j$ is a cubic, and $q_j(0)=0$, $\nabla q_j(0)=0$ because $\phi$ and $\nabla\phi$ vanish at $z$ and $L_j(z)=1$. On the boundary ray through $p_0$, $N_0^2$ divides $q_1$, hence each homogeneous part. For the gluing, $L_{j+1}-L_j$ is linear and vanishes at $p_j$, so it is a multiple of $N_j$. Modulo $N_j^2$, $L_{j+1}^2\equiv L_j^2+2\beta_jN_jL_j$, so
\[
\phi_j-\phi_{j+1}\equiv L_j^2(q_j-q_{j+1})-2\beta_jN_jL_jq_{j+1}\pmod{N_j^2}.
\]
This vanishes iff $q_j-q_{j+1}=N_jg$ and $L_jg-2\beta_jq_{j+1}=0$ on the ray. Write $g=G_1+G_2$ (degrees 1, 2). On the ray, $x=tp_j$, $L_j=1-t$, and the second condition reads $tG_1(p_j)+t^2\bigl(G_2(p_j)-G_1(p_j)-2\beta_jQ_{j+1}(p_j)\bigr)-t^3\bigl(G_2(p_j)+2\beta_jC_{j+1}(p_j)\bigr)=0$. Hence $G_1(p_j)=0$, which with $N_j\mid Q_j-Q_{j+1}$ is (G1). $N_j\mid C_j-C_{j+1}$ is (G2). $G_2(p_j)=2\beta_jQ_j(p_j)$ is (G3), because $\nabla(N_jG_2)(p_j)=G_2(p_j)\nabla N_j$ and $Q_j(p_j)=Q_{j+1}(p_j)$ by (G1). Finally $G_2(p_j)=-2\beta_jC_{j+1}(p_j)$ gives (G4) with (G2). Every step reverses.
\end{proof}

\paragraph{The witness.} From now on $F=(0;p_0,p_1,p_2,p_3)$ is a fan of three triangles with $D_{01},D_{12},D_{23}>0$. Put
\[
A:=D_{13}D_{23},\qquad A_3:=D_{01}D_{02}.
\]
If $p_i=r_i(\cos\vartheta_i,\sin\vartheta_i)$ with $\vartheta_0=0$, $\vartheta_j=\alpha_1+\dots+\alpha_j$, then
\begin{equation}\label{m3:signs}
A=r_1r_2r_3^2\sin(\alpha_2+\alpha_3)\sin\alpha_3,\qquad A_3=r_0^2r_1r_2\sin\alpha_1\sin(\alpha_1+\alpha_2),
\end{equation}
so $A,A_3>0$ whenever $\alpha_1+\alpha_2<\pi$ and $\alpha_2+\alpha_3<\pi$.

\begin{definition}[Witness]\label{m3:wit}
Let $\ell_1$ be the linear form with $\ell_1(p_0)=-4A$, $\ell_1(p_1)=-A$, and $\ell_3$ the one with $\ell_3(p_3)=-4A_3$, $\ell_3(p_2)=-A_3$. Let
\[
Q_1=AN_0^2,\qquad Q_2=AN_0^2-\frac{D_{02}D_{03}D_{23}}{D_{12}}N_1^2,\qquad Q_3=A_3N_3^2,\qquad C_1=N_0^2\ell_1,\qquad C_3=N_3^2\ell_3,
\]
and let $C_2$ be the cubic form with polar values
\begin{gather*}
C_2(p_1,p_1,p_1)=-AD_{01}^2,\qquad C_2(p_2,p_2,p_2)=-A_3D_{23}^2,\\
C_2(p_1,p_1,p_2)=\tfrac{A D_{01}}{3}(6D_{12}-5D_{02}+2D_{01}),\qquad C_2(p_1,p_2,p_2)=\tfrac{A_3D_{23}}{3}(6D_{12}-5D_{13}+2D_{23}).
\end{gather*}
Set $\phi_j:=L_j^2(Q_j+C_j)$ on $T_j$ and $w:=\curl\phi$. In the barycentric coordinates $(\lambda_0,\lambda_1,\lambda_2)$ of $T_j$ ($\lambda_0=L_j$ at $z$, $\lambda_1$ at $p_{j-1}$, $\lambda_2$ at $p_j$),
\begin{align*}
\phi_1&=AD_{01}^2\,\lambda_0^2\lambda_2^2(\lambda_0-3\lambda_1),\qquad \phi_3=A_3D_{23}^2\,\lambda_0^2\lambda_1^2(\lambda_0-3\lambda_2),\\
\phi_2&=\lambda_0^2\Bigl[AD_{01}^2\lambda_1^2(\lambda_0+\lambda_2)+A_3D_{23}^2\lambda_2^2(\lambda_0+\lambda_1)+2AD_{01}D_{02}\,\lambda_1\lambda_2(\lambda_0+\lambda_1+\lambda_2)\\
&\qquad\ +AD_{01}(6D_{12}-5D_{02}+2D_{01})\,\lambda_1^2\lambda_2+A_3D_{23}(6D_{12}-5D_{13}+2D_{23})\,\lambda_1\lambda_2^2\Bigr].
\end{align*}
\end{definition}

\begin{lemma}[The witness is admissible]\label{m3:witlem}
$\phi\in C^1(F)$ is piecewise quintic with $\phi=|\nabla\phi|=0$ on $\partial F$. On both outer rays $\int\partial_{nn}\phi=0$, and
\[
\int_{[0,p_0]}s^2\,\partial_{nn}\phi\,ds=\frac{|p_0|^5A}{30},\qquad \int_{[0,p_3]}s^2\,\partial_{nn}\phi\,ds=\frac{|p_3|^5A_3}{30}.
\]
Hence $w$ is a continuous, divergence-free, piecewise $P_4$ field vanishing on $\partial F$. Extended by zero it lies in $Y_h^1$, both when $F$ is the whole star ($m=3$) and when $F=T_1\cup T_2\cup T_3$ with $m\ge4$. In the first case $|\mathcal M(w)|=\bigl||p_0|^5A+|p_3|^5A_3\bigr|/60$. In the second, $|\mathcal M(w)|=|p_0|^5|A|/60$.
\end{lemma}

\begin{proof}
\emph{Boundary.} The factor $L_j^2$ gives first-order vanishing on the outer edges, and $N_0^2\mid q_1$, $N_3^2\mid q_3$ on the outer rays.

\emph{Ray through $p_1$.} (G1) holds by construction, and (G2)$+$(G4) read $C_1(p_1)=C_2(p_1)=-Q_1(p_1)=-AD_{01}^2$. This holds since $\ell_1(p_1)=-A$ and $N_0(p_1)=D_{01}$. (G3) is an identity between two vectors, so we test it with $p_1$ and with $p_2$.
\begin{itemize}
\item The $p_1$-component. By Euler's identity it reads $3(C_1-C_2)(p_1)=2\beta_1Q_1(p_1)N_1(p_1)$. Both sides vanish, by (G2) and because $N_1(p_1)=0$.
\item The $p_2$-component. It reads $3C_1(p_1,p_1,p_2)-3C_2(p_1,p_1,p_2)=2\beta_1AD_{01}^2D_{12}$. Here $3C_1(p_1,p_1,p_2)=D_{01}^2\ell_1(p_2)+2D_{01}D_{02}\ell_1(p_1)=AD_{01}(4D_{12}-3D_{02})$, using $p_2=(D_{02}p_1-D_{12}p_0)/D_{01}$. Also $\beta_1=(D_{02}-D_{12}-D_{01})/(D_{01}D_{12})$, obtained by evaluating $L_2-L_1$ at $p_2$.
\end{itemize}
Substituting gives exactly the prescribed $C_2(p_1,p_1,p_2)$.

\emph{Ray through $p_2$.} This is symmetric. (G1) is the identity
\begin{equation}\label{m3:plucker}
AN_0^2-\frac{D_{02}D_{03}D_{23}}{D_{12}}N_1^2+\frac{D_{01}D_{03}D_{13}}{D_{12}}N_2^2-A_3N_3^2\equiv0,
\end{equation}
a three-term Pl\"ucker relation among the squares of four linear forms in two variables. In angular form it is $\sum_j N_j^2/\bigl(r_j^2\prod_{k\ne j}\sin(\vartheta_j-\vartheta_k)\bigr)=0$.

\emph{Means and moments.} On $[0,p_0]$ we have $\phi_1=L_1^2N_0^2(A+\ell_1)$ and $N_0=0$. With $n$ the unit normal, $\partial_{nn}\phi=2|\nabla N_0|^2L_1^2(A+\ell_1)=2|p_0|^2L_1^2(A+\ell_1)$. At $x=tp_0$ this is $2|p_0|^2A\,(1-t)^2(1-4t)$, because $\ell_1(tp_0)=-4At$. Now $\int_0^1(1-t)^2(1-4t)\,dt=0$ and $\int_0^1(t-\frac12)^2(1-t)^2(1-4t)\,dt=\frac1{60}$, and $ds=|p_0|\,dt$, $s=|p_0|(t-\frac12)$. The ray $[0,p_3]$ is symmetric. The formula for $\mathcal M$ follows from \eqref{m3:dnv}. By (G), the only chords meeting $F$ are $[0,p_0]$ and, when $m=3$, $[0,p_3]$; on all others the trace of $\dn w$ from $T_e$ vanishes.

All identities in this proof are verified as identities of rational functions of the free coordinates of $p_1,p_2,p_3$ (with $p_0=(1,0)$, which is no loss by similarity) in \texttt{m3\_certify.py symbolic}. The barycentric form is checked there as well.
\end{proof}

\emph{Reflection.} Reflect the fan in a line through $z$ and reverse the order of the rays. Then $D_{ij}\mapsto D_{3-j,3-i}$, so $A\leftrightarrow A_3$, $N_1\leftrightarrow N_3$ and $N_2$ is unchanged in Lemma~\ref{m3:norm} below. Hence the witness data are reflection invariant; this is checked exactly in the log, including a fan with $\beta_1=0$.

\begin{lemma}[The norm]\label{m3:norm}
For a polynomial $f$ in the barycentrics of a triangle $T$ with angles $g_0,g_1,g_2$,
\[
\norm{D^2f}_{L^2(T)}^2=\frac{P_f(\cot g_0,\cot g_1,\cot g_2)}{2|T|}
\]
with $P_f$ a quadratic form with rational coefficients. This follows from $2|T|\nabla\lambda_a\cdot\nabla\lambda_a=\cot g_b+\cot g_c$, $2|T|\nabla\lambda_a\cdot\nabla\lambda_b=-\cot g_c$ and $\int_T\lambda^\alpha=2|T|\alpha!/(|\alpha|+2)!$. Consequently
\begin{gather*}
\norm{\nabla w}^2=N_1+N_2+N_3,\qquad N_1=A^2D_{01}^3P_1(\cot T_1),\qquad N_3=A_3^2D_{23}^3P_3(\cot T_3),\\ N_2=\frac{k^{\mathsf T}K(\cot T_2)\,k}{D_{12}},\\
k=\bigl(AD_{01}^2,\ A_3D_{23}^2,\ 2AD_{01}D_{02},\ AD_{01}(6D_{12}-5D_{02}+2D_{01}),\ A_3D_{23}(6D_{12}-5D_{13}+2D_{23})\bigr),
\end{gather*}
where $\cot T_j$ lists the cotangents at $(z,p_{j-1},p_j)$. The forms are
\[
P_1=\tfrac{19}{70}k_0^2+\tfrac1{10}k_0k_1+\tfrac{17}{35}k_0k_2+\tfrac3{70}k_1^2+\tfrac17k_1k_2+\tfrac3{10}k_2^2,
\]
$P_3=P_1$ with $k_1\leftrightarrow k_2$, and $K$ is the $5\times5$ Gram matrix of the five bracketed terms of $\phi_2$ printed in the log. For a star with $m=3$ put
\[
\kappa_w:=\frac{|p_0|^5A+|p_3|^5A_3}{60\,\max(|p_0|,|p_3|)^2\sqrt{N_1+N_2+N_3}},
\]
and for $m\ge4$ put $\kappa_w:=|p_0|^5A/\bigl(60|p_0|^2\sqrt{N_1+N_2+N_3}\bigr)$. Then $w/\norm{\nabla w}$, with sign chosen, satisfies (M3)(iii) with constant $\kappa_w$.
\end{lemma}

The formula is checked in the log against direct exact integration of $|D^2\phi|^2$ on random rational fans.

\begin{lemma}[Uniqueness for $m=3$]\label{m3:dim}
Let $m=3$, $\beta_1\beta_2\ne0$, and suppose no two of the rays through $p_0,\dots,p_3$ lie on one line, except possibly those through $p_0$ and $p_3$. Then $Y^1(\omega_z):=\{v\in Y_h^1:\ \operatorname{supp}v\subset\omega_z\}$ is spanned by $w$. In particular the star constant is $\kappa_z=\kappa_w$ (normalised by $|e_z|$), and $\kappa_z=0$ iff $|p_0|^5A+|p_3|^5A_3=0$.
\end{lemma}

\begin{proof}
By Lemma~\ref{m3:glue}, the quadratic parts satisfy $aN_0^2+c_1N_1^2+c_2N_2^2-a'N_3^2=0$. The squares of three pairwise independent linear forms in two variables are linearly independent. So this relation is unique up to scale, and it is \eqref{m3:plucker} (if $N_3\parallel N_0$, it forces $c_1=c_2=0$, $a=a'$). If the quadratic part vanishes, (G4) with $\beta_j\ne0$ gives $C_{j+1}(p_j)=C_j(p_j)=0$, so each $C_j$ is divisible by the forms of both of its rays: $C_1=\alpha N_0^2N_1$, $C_2=N_1N_2\mu$, $C_3=\alpha_3N_3^2N_2$. The means on the two chords are $\propto\alpha N_1(p_0)$ and $\alpha_3N_2(p_3)$, so $\alpha=\alpha_3=0$. Then (G3) at $p_1$ and $p_2$ gives $\mu(p_1)=\mu(p_2)=0$, so $\mu=0$. Hence $\dim Y^1(\omega_z)\le1$, and $w\ne0$.
\end{proof}

The hypotheses of Lemma~\ref{m3:dim} are checked exactly for every star of the production meshes (Section~\ref{m3:sec:mesh}). The exact dimension counts in the log agree: $\dim Y(\omega_z)=2m-3$ and $\dim Y^1(\omega_z)=2m-5$. Two cases behave differently: when $\beta_1=0$ the dimension jumps to $2$, and the witness is still admissible; a fan of one or two triangles carries nothing.

\begin{remark}[The hypothesis $\alpha_1+\alpha_2,\alpha_2+\alpha_3<\pi$ is needed]\label{m3:neg}
If $\alpha_1+\alpha_2>\pi$, then $A_3<0<A$ and the two terms of $\mathcal M(w)$ compete. The log exhibits a segment of $m=3$ stars on which $|p_0|^5A+|p_3|^5A_3$ changes sign, with $\beta_1\beta_2\ne0$. At the zero, by Lemma~\ref{m3:dim}, no field supported in the star satisfies (M3). Such stars have $\Theta_z>\pi+\alpha_3$, which (A$_\delta$) excludes. The coarsest production mesh $N=8$ has $\alpha_1+\alpha_2>\pi$ at every vertex ($\Theta_z=225^\circ$), but there the sum is still positive: $\kappa_z\ge0.00287$.
\end{remark}

\paragraph{Qualitative statement and compactness.}
\begin{definition}[Families]\label{m3:fam}
$F_3(\theta_{\min},\delta)$ is the set of three-triangle fans with all nine angles $\ge\theta_{\min}$ and $|\alpha_1+\alpha_2+\alpha_3-\pi|\le\delta$. $F_s(\theta_{\min},\delta)$ is the set with all nine angles $\ge\theta_{\min}$ and $\alpha_1+\alpha_2+\alpha_3\le\pi+\delta-\theta_{\min}$. Both are taken modulo similarity, normalised by $|p_0|=1$ and parametrised by $(\alpha_j,\beta_j)_{j=1}^3$. They are compact.
\end{definition}

Under (M0$_\theta$), (M1$'$), (A$_\delta$) the star is in $F_3$ if $m=3$. If $m\ge4$, then $T_1\cup T_2\cup T_3$ (numbered from $e_z$, after a reflection if necessary) is in $F_s$, since the remaining $m-3\ge1$ angles at $z$ are $\ge\theta_{\min}$. On both families every pair sum $\alpha_j+\alpha_{j+1}$ lies in $[2\theta_{\min},\pi+\delta-\theta_{\min}]\subset(0,\pi)$, so $A,A_3>0$ by \eqref{m3:signs}. Hence $\mathcal M(w)\ne0$, $w\ne0$ and $\kappa_w>0$. $\kappa_w$ is a continuous function of $(\alpha_j,\beta_j)$, since $D_{01},D_{12},D_{23}$ stay positive. So
\[
\kappa_\ast(\theta_{\min},\delta):=\min\Bigl(\min_{F_3}\kappa_w,\ \min_{F_s}\kappa_w\Bigr)>0 .
\]
By the reflection invariance, $\min_{F_3}\kappa_w=\min_{F_3\cap\{|p_3|\le|p_0|\}}\kappa_w$, and on the latter set $\kappa_w$ is normalised by $|p_0|$.

\emph{Proof of Theorem~\ref{m3:thm}, main part.} Take $v_z=\pm w/\norm{\nabla w}$ on $\omega_z$. For $m=3$, $\kappa_w$ is normalised by the longer chord, so (iii) holds whichever chord is $e_z$. (i) holds by construction, and (iii) by Lemmas~\ref{m3:witlem} and \ref{m3:norm}. For (ii), $\omega_z$ lies within $\max_jr_j\le(\sin\theta_{\min})^{-3}|e_z|$ of $z$ (law of sines, three triangles), and a triangle lies in at most three patches, one per vertex. \qed

\paragraph{A decomposed bound.}
For the certificates and for (a) we bound $\kappa_w$ from below by an expression whose factors each involve only one or two triangles. Normalise $|p_0|=1$. Let $q_j:=r_j/r_{j-1}=\sin\beta_j/\sin\gamma_j$, $s_j:=\sin\alpha_j$, $s_{12}:=\sin(\alpha_1+\alpha_2)$, $s_{23}:=\sin(\alpha_2+\alpha_3)$ and $\rho:=s_1s_{12}/(s_{23}s_3)$. Write $\norm{u}_K^2:=u^{\mathsf T}K(\cot T_2)u$.

\begin{lemma}[Decomposition]\label{m3:dec}
On $F_3\cap\{r_3\le1\}$ and on $F_s$,
\[
\kappa_w\ \ge\ \frac{b}{\sqrt{3600\bigl(F_1+(\sqrt U+\sqrt W)^2+\Phi_3\bigr)}},
\]
with $b=1+r_3^3\rho$ on $F_3$ and $b=1$ on $F_s$, where
\begin{align*}
F_1&=(q_1s_1)^3P_1(\cot T_1),\qquad U=\frac{q_1^2}{q_2s_2}\norm{\tilde u}_K^2,\qquad W=\frac{\rho^2q_1^2q_2}{s_2}\norm{\tilde w}_K^2,\\
\Phi_3&=\frac{\rho^2(q_1q_2)^2s_3^3}{q_3}P_3(\cot T_3),\qquad \tilde u=\bigl(s_1^2,\,0,\,2s_1q_2s_{12},\,s_1(6q_1q_2s_2-5q_2s_{12}+2s_1),\,0\bigr),\\
\tilde w&=\bigl(0,\,q_2s_3^2,\,0,\,0,\,s_3(6s_2/q_3-5s_{23}+2q_2s_3)\bigr).
\end{align*}
\end{lemma}

\begin{proof}
By \eqref{m3:signs}, $A=r_1r_2\,X$ with $X=r_3^2s_{23}s_3$ and $A_3=r_1r_2\,X\rho/r_3^2$. So $|p_0|^5A+|p_3|^5A_3=r_1r_2X(1+r_3^3\rho)$, and the common factor $(r_1r_2)^2$ cancels between $\mathcal M^2$ and $N_1+N_2+N_3$. Dividing the latter by $(r_1r_2X)^2$ gives:
\begin{itemize}
\item $N_1$ gives $F_1$ ($D_{01}=q_1s_1$) and $N_3$ gives $\Phi_3$.
\item $N_2=\norm{k}_K^2/D_{12}$. Substituting $D_{01}=q_1s_1$, $D_{12}=q_1^2q_2s_2$, $D_{02}=q_1q_2s_{12}$, $D_{13}=r_1r_3s_{23}$, $D_{23}=r_2r_3s_3$ gives $k/(r_1r_2X)=q_1^2(\tilde u+\rho\,q_2\tilde w)$. Minkowski's inequality in the seminorm $\norm\cdot_K$ ($K$ is a Gram matrix) then gives $\sqrt{N_2}/(r_1r_2X)\le\sqrt U+\sqrt W$.
\end{itemize}
On $F_3\cap\{r_3\le1\}$ the normalising chord is $|p_0|=1$.
\end{proof}

\begin{corollary}[Elementary bound]\label{m3:el}
Let $\sigma=\sin\theta_{\min}$, $c=\cot\theta_{\min}$, $\sigma_\ast=\min(\sin2\theta_{\min},\sin(\theta_{\min}-\delta))$, $Q=1/\sigma$, $R=1/(\sigma\sigma_\ast)$. For a polynomial $p$ let $\norm p_1$ be the sum of the absolute values of its coefficients, and $\kappa_{ij}:=\norm{K_{ij}}_1$. Then $\kappa_\ast(\theta_{\min},\delta)\ge\kappa_{\rm el}:=\bigl(3600(\bar F_1+(\sqrt{\bar U}+\sqrt{\bar W})^2+\bar\Phi_3)\bigr)^{-1/2}$, with
\begin{gather*}
\bar F_1=Q^3\norm{P_1}_1c^2,\qquad \bar U=Q^4\,c^2\sum_{i,j}\bar u_i\kappa_{ij}\bar u_j,\\ \bar W=R^2Q^4c^2\sum_{i,j}\bar w_i\kappa_{ij}\bar w_j,\qquad \bar\Phi_3=R^2Q^5\norm{P_3}_1c^2,\\
\bar u=(1,0,2Q,6Q^2+5Q+2,0),\qquad \bar w=(0,Q,0,0,8Q+5).
\end{gather*}
For instance $\kappa_{\rm el}(30^\circ,10^\circ)=5.1\cdot10^{-5}$ and $\kappa_{\rm el}(20^\circ,15^\circ)=2.4\cdot10^{-6}$ (log).
\end{corollary}

\begin{proof}
Every angle lies in $[\theta_{\min},\pi-2\theta_{\min}]$, so $\sin\ge\sigma$ and $|\cot|\le c$. Hence $q_j\in[\sigma,1/\sigma]$, and $\rho\le R$ because the pair sums lie in $[2\theta_{\min},\pi+\delta-\theta_{\min}]$. Also $|P(k)|\le\norm P_1c^2$ for the quadratic forms involved. Insert these into Lemma~\ref{m3:dec} with $b\ge1$. This uses only $A_3\ge0$, not $r_3\le1$, so it bounds $\kappa_w$ normalised by $|p_0|$ on all of $F_3$ and $F_s$; for $F_3$ use reflection when $e_z=[z,p_3]$.
\end{proof}

The elementary bound is valid for all $\theta_{\min}\le\pi/3$ and $\delta<\theta_{\min}$. It is pessimistic: about $25$ times below the certified value at $(30^\circ,10^\circ)$ and about $125$ times below at $(20^\circ,15^\circ)$.

\paragraph{Interval certificates.}\phantomsection\label{m3:sec:cert}
For fixed $(\theta_{\min},\delta)$ we certify $\kappa_w\ge\kappa_0$ on $F_3\cap\{r_3\le1\}$ and on $F_s$ by branch and bound over boxes in $(\alpha_1,\alpha_2,S\text{ or }\alpha_3,\beta_1,\beta_2,\beta_3)$, where $S=\alpha_1+\alpha_2+\alpha_3$ for $F_3$. Each box is first \emph{contracted}: its upper and lower ends are moved inwards using the linear constraints of the family ($\alpha_j+\beta_j\le\pi-\theta_{\min}$; $\theta_{\min}\le S-\alpha_1-\alpha_2\le\pi-2\theta_{\min}$ for $F_3$; $\sum\alpha_j\le\pi+\delta-\theta_{\min}$ for $F_s$). The new ends are rounded outward, so the contracted box still contains every point of the family that the original box contained. A box is then
\begin{itemize}
\item \emph{discarded} if it provably misses the family: it is empty after contraction, or some $\gamma_j<\theta_{\min}$ on the whole box, or $\alpha_3<\theta_{\min}$, or $S$ out of range;
\item \emph{covered by symmetry} ($F_3$ only) if $r_3>1$ on the whole box, since its reflection has $r_3<1$;
\item \emph{certified} if an interval evaluation of Lemma~\ref{m3:dec} gives $b^2\ge3600\kappa_0^2\,(F_1+(\sqrt U+\sqrt W)^2+\Phi_3)$ on the whole box;
\item \emph{bisected} along its widest side otherwise.
\end{itemize}
Rigour rests on the following.
\begin{itemize}
\item All arithmetic is in IEEE double precision with every result widened outward by one ulp; basic operations are correctly rounded.
\item $\sin$ and $\cos$ at box endpoints are Taylor polynomials of degree $41$ evaluated in interval arithmetic, plus the Lagrange remainder. They are extended to intervals by monotonicity, including the extrema at $\pi/2$ and $3\pi/2$.
\item $\cot$ is enclosed by monotonicity on $(0,\pi)$. $q_j$ is enclosed by its exact corner values when $\gamma_j$ stays on one side of $\pi/2$, where it is monotone in $\alpha_j$ and $\beta_j$.
\item $\pi$ is enclosed by the two doubles adjacent to it.
\item The rational coefficients of $P_1,P_3,K$ come from exact \texttt{SymPy} derivations and are enclosed in outward-rounded doubles.
\end{itemize}
Results (log, section 5):
\begin{center}
\small
\begin{tabular}{llll}
\toprule
family & certified $\kappa_w\ge$ & boxes (cert./disc./sym./total) & float min.\\
\midrule
$F_3(30^\circ,10^\circ)$ & $0.0020$ & $203150/154/6503/419613$ & $0.002558$\\
$F_s(30^\circ,10^\circ)$ & $0.0012$ & $209555/0/0/419109$ & $0.001400$\\
$F_3(20^\circ,15^\circ)$ & $0.0005$ & $1326711/9699/38333/2749485$ & $0.000880$\\
$F_s(20^\circ,15^\circ)$ & $0.0003$ & $452009/0/0/904017$ & $0.000395$\\
\bottomrule
\end{tabular}
\end{center}
As a negative control, the same code with $\kappa_0$ set $5\%$ above the float minimum of the bound on $F_s(30^\circ,10^\circ)$ does not certify. It is refuted when a box of width $<10^{-9}$ cannot be certified.

\paragraph{The production meshes.}\phantomsection\label{m3:sec:mesh}
For the meshes of Section~\ref{sec:numerics} we evaluate $\kappa_w$ exactly. The double-precision vertex coordinates are read as exact rationals, $D_{ij}$ and all cotangents are rational, and $|p_0|$, $|p_3|$ are enclosed by rational square roots. We check exactly that $D_{01},D_{12},D_{23}>0$, $D_{02}D_{13}\ne0$ and $\beta_1\beta_2\ne0$ at every boundary vertex. By Lemma~\ref{m3:dim} the resulting numbers are therefore the exact star constants. They reproduce the floating-point Table~\ref{lb:tab:stars} to all printed digits, for example $\min_z\kappa_z=0.002879$ ($N=8$), $0.005551$ ($N=16$), $0.003663$ ($N=64$), $0.002876$ ($N=1024$). The minimum over both choices of $e_z$ is $\ge0.00287$ on all levels.

\cert{code/m3\_certify.py all; notes/v6/m3/recert.py}{logs/m3\_certify.log, notes/v6/m3/recert\_30\_10.log}
The interval code is hand-written outward-rounded double arithmetic, not a verified library; Theorem~\ref{m3:thm} itself (positivity of $\kappa_\ast$) does not depend on it.

\paragraph{What is not covered.} The certified numerical constants are for $(\theta_{\min},\delta)$ as listed; for other values Theorem~\ref{m3:thm} gives $\kappa_\ast>0$ by compactness and the explicit, pessimistic $\kappa_{\rm el}$. The theorem needs $\delta<\theta_{\min}$, which (S) guarantees; coarse meshes with $\Theta_z-\pi\ge\min(\alpha_1,\alpha_3)$ need the star-by-star check of Section~\ref{m3:sec:mesh}. By Remark~\ref{m3:neg} (M3) can fail on such stars for star-supported fields. For $m\ge4$ we use only three of the $m$ triangles, which is enough but not optimal.


\subsubsection{Three lemmas}

\begin{lemma}\label{lb:normal}
If $v\in Y_h$ and $e\in\EG$, then $\dn v\cdot n_e=0$ on $e$.
\end{lemma}

\begin{proof}
$v|_{T_e}$ is a polynomial that vanishes on $e$, so $\dt v=0$ on $e$. Then $0=\div v=\dt(v\cdot\tau_e)+\dn(v\cdot n_e)$ on $e$.
\end{proof}

\begin{lemma}[The error identity on $Y_h$]\label{lb:identity}
Let $u_h$ be the velocity of $\CNS$ ($\gamma\ge\gamma_2$), or of $\GS$ ($\gamma\ge\gamma_0$, any $G$), and let $e_u:=\ut-u_h$. Then for all $v\in Y_h$,
\[
a_h(e_u,v)-\ip{e_u}{\dn v}=-\ip{\ut}{\dn v}-(F,v).
\]
\end{lemma}

\begin{proof}
For $\CNS$ use \eqref{eq:cnserr}: $B^\ast(e_p,v)=-(e_p,\div v)+\ip{e_p-\pbG(e_p)}{v\cdot n}=0$, because $\div v=0$ and $v=0$ on $\Gh$. For $\GS$ use Corollary~\ref{cor:GSerr}: $\Rc(v)=-\ip{\pt-c}{v\cdot n}=0$. In both cases $N_h(e_u,v)=\Gc(v)$. Since $v=0$ on $\Gh$, $N_h(e_u,v)=a_h(e_u,v)-\ip{e_u}{\dn v}$, and in $\Gc(v)$ the transposed and penalty terms vanish, leaving $\Gc(v)=-\ip{\ut}{\dn v}-(F,v)$.
\end{proof}

Neither the pressure error nor the penalty enters this identity. It is simply the discrete momentum equation $a_h(u_h,v)-\ip{u_h}{\dn v}=(\ft,v)$ on $Y_h$, which $\ut$ violates by $-\ip{\ut}{\dn v}$ (Nitsche's symmetric term applied to the nonzero trace of $\ut$ on the chords) and by $(F,v)$.

\begin{lemma}[Trace control on $Y_h^1$]\label{lb:poinc}
Let $C_{PT}$ be the constant in $\norm{w-\bar w_T}_{L^2(e)}\le C_{PT}|e|^{1/2}\norm{\nabla w}_{T}$ for $e\subset\partial T$, where $\bar w_T$ is the mean of $w$ over $T$; it depends only on shape regularity. Then for $w\in H^1(\Oh)^2$ and $v\in Y_h^1$,
\[
|\ip{w}{\dn v}|\le C_{PT}C_I\norm{\nabla w}\norm{\nabla v}.
\]
\end{lemma}

\begin{proof}
On each $e$, $\int_e\dn v=0$, so $\int_e w\cdot\dn v=\int_e(w-\bar w_{T_e})\cdot\dn v$. Hence
\[
|\ip{w}{\dn v}|\le\sum_e C_{PT}\norm{\nabla w}_{T_e}\,|e|^{1/2}\norm{\dn v}_{L^2(e)}\le C_{PT}C_I\sum_e\norm{\nabla w}_{T_e}\norm{\nabla v}_{T_e},
\]
using the local form of Lemma~\ref{lem:inverse}. Finish with Cauchy--Schwarz; the $T_e$ are distinct.
\end{proof}

\begin{corollary}[Abstract lower bound]\label{lb:abstract}
With $\Lambda_0:=2+C_{PT}C_I$ and $e_u=\ut-u_h$ as in Lemma~\ref{lb:identity},
\[
\Lambda_0\norm{\nabla e_u}\ \ge\ \sup_{0\ne v\in Y_h^1}\frac{|\ip{\ut}{\dn v}+(F,v)|}{\norm{\nabla v}} .
\]
\end{corollary}

\begin{proof}
$|a_h(e_u,v)|\le2\norm{\nabla e_u}\norm{\nabla v}$; apply Lemmas~\ref{lb:identity} and~\ref{lb:poinc}.
\end{proof}

\begin{lemma}[Leading term]\label{lb:expand}
For $v\in Y_h^1$,
\[
\Bigl|\ip{\ut}{\dn v}-\ell_W(v)\Bigr|\le C\hG^{5/2}\norm{\nabla v},\qquad \ell_W(v):=-\frac12\sum_{e\in\EG}W_e\int_e s^2\,\dn v\cdot\tau_e\,ds,
\]
and $|(F,v)|\le C\hG^2\norm{\nabla v}$.
\end{lemma}

\begin{proof}
\begin{enumerate}
\item \emph{Geometry.} Let $x\in e$ and let $x^\ast=x/|x|$ be its nearest point on $\Gamma$. Then $x=x^\ast+d\,n(x^\ast)$ with $d=1-|x|$. Put $\ell=|e|$. Since $|x|^2=1-\ell^2/4+s^2$,
\[
d=\tfrac12\bigl(\tfrac{\ell^2}4-s^2\bigr)+O(\ell^4).
\]
\item \emph{Taylor expansion.} Using $\ut(x^\ast)=0$ and Lemma~\ref{lem:wall},
\[
\ut(x)=d\,W(x^\ast)\tau(x^\ast)+O(d^2).
\]
Moreover $\tau(x^\ast)\cdot\tau_e=1+O(s^2)$ and $W(x^\ast)=W_e+O(|s|)$. Hence
\[
\ut\cdot\tau_e=\tfrac12\bigl(\tfrac{\ell^2}{4}-s^2\bigr)W_e+O(\hG^3)\quad\text{on }e .
\]
\item \emph{Pairing on one edge.} By Lemma~\ref{lb:normal}, $\ut\cdot\dn v=(\ut\cdot\tau_e)(\dn v\cdot\tau_e)$ on $e$. The constant part $\frac{\ell^2}{8}W_e$ integrates to zero against $\dn v\cdot\tau_e$, whose mean vanishes because $v\in Y_h^1$. Therefore
\[
\int_e\ut\cdot\dn v=-\tfrac12W_e\int_es^2\,\dn v\cdot\tau_e+O(\hG^3)\norm{\dn v}_{L^1(e)} .
\]
\item \emph{Summation.} $\sum_e\norm{\dn v}_{L^1(e)}\le\sum_e|e|^{1/2}\norm{\dn v}_{L^2(e)}\le C_I(\#\EG)^{1/2}\norm{\nabla v}\le C\hG^{-1/2}\norm{\nabla v}$.
\item \emph{The $F$ term.} This is Lemma~\ref{lem:ext}.\qedhere
\end{enumerate}
\end{proof}

\begin{lemma}[Assembling the test field]\label{lb:assemble}
Assume (M3), $W\not\equiv0$, and $\hG\le h_1$, where $h_1$ depends on $\norm{W}_{W^{1,\infty}(\Gamma)}/\norm{W}_{L^2(\Gamma)}$, on (M0), on $C_I$, and on the constants $\kappa_0$, $C_s$, $K_{\mathrm{ov}}$ of (M3). Then $v:=\sum_zW(z)\,v_z\in Y_h^1$ satisfies
\[
\ell_W(v)\ \ge\ \tfrac12\kappa_0c_0^{5/2}K_{\mathrm{ov}}^{-1/2}\,\norm{W}_{L^2(\Gamma)}\,\hG^{3/2}\,\norm{\nabla v}.
\]
\end{lemma}

\begin{proof}
\begin{enumerate}
\item \emph{Norm.} By bounded overlap, $\norm{\nabla v}^2\le K_{\mathrm{ov}}\sum_zW(z)^2$.
\item \emph{Splitting $\ell_W$.} Write $m_{z,e}:=-\frac12\int_es^2\,\dn v_z\cdot\tau_e$, so that $\sum_em_{z,e}=\mathcal M(v_z)$. Then $|m_{z,e}|\le\frac{\hG^2}{8}C_I$, and $m_{z,e}\ne0$ only for the boundedly many edges $e$ in $\omega_z$. For those, $|W_e-W(z)|\le C\hG\norm{W}_{W^{1,\infty}}$. Hence
\[
\ell_W(v)=\sum_zW(z)\sum_eW_e\,m_{z,e}\ \ge\ \sum_zW(z)^2\mathcal M(v_z)-C\hG^3\norm{W}_{W^{1,\infty}}\sum_z|W(z)| .
\]
\item \emph{Lower bound.} By (M3)(iii) and (M0), $\mathcal M(v_z)\ge\kappa_0c_0^2\hG^2$. Moreover $\sum_z|W(z)|\le\#\EG\,\norm{W}_\infty\le C\hG^{-1}\norm{W}_\infty$.
\item \emph{Riemann sum.} Let $\Sigma_W:=\sum_zW(z)^2|e_z|$. Since $z\mapsto e_z$ is a bijection onto the edges, this is a Riemann sum for $\int_\Gamma W^2$, and $|\Sigma_W-\norm{W}^2_{L^2(\Gamma)}|\le C\hG\norm{W}^2_{W^{1,\infty}}$; and $\sum_zW(z)^2\in[\Sigma_W/\hG,\ \Sigma_W/(c_0\hG)]$.
\item Combining,
\[
\ell_W(v)\ge\kappa_0c_0^2\hG\,\Sigma_W-C\hG^2\norm{W}^2_{W^{1,\infty}},\qquad \norm{\nabla v}\le(K_{\mathrm{ov}}\Sigma_W/(c_0\hG))^{1/2}.
\]
For $\hG\le h_1$ we have $\Sigma_W\ge\frac34\norm{W}^2_{L^2(\Gamma)}$, and the negative term is at most a quarter of the positive one.\qedhere
\end{enumerate}
\end{proof}

\subsubsection{The theorem}

\begin{theorem}[$H^1$ lower bound]\label{lb:thm}
Assume (M0)--(M2), (D) and (M3). By Theorem~\ref{m3:thm}, (M3) holds for every $k\ge4$ when $\hG\le2\sin(\theta_{\min}/4)$, with $\kappa_0=\kappa_\ast(\theta_{\min},\theta_{\min}/2)$, $K_{\mathrm{ov}}=3$ and $C_s=(\sin\theta_{\min})^{-3}$. Let $W=\dn u\cdot\tau\not\equiv0$ on $\Gamma$. Let $u_h$ be the velocity of $\CNS$ with $\gamma\ge\gamma_2$, or of $\GS$ with $\gamma\ge\gamma_0$ and any $G$. Then, for $\hG\le h_1$,
\[
\norm{\nabla(\ut-u_h)}_{\Oh}\ \ge\ c_{\mathrm{LB}}\,\kappa_0\,\norm{W}_{L^2(\Gamma)}\,\hG^{3/2}-C\hG^{2},\qquad c_{\mathrm{LB}}:=\frac{c_0^{5/2}}{2\Lambda_0K_{\mathrm{ov}}^{1/2}} .
\]
Here $c_{\mathrm{LB}}$ depends on $c_0$, on $K_{\mathrm{ov}}$ from (M3), and on $C_{PT}$ and $C_I$ (hence on (M0)--(M2)); $h_1$ is as in Lemma~\ref{lb:assemble}; and $C$ depends on $\norm{\ut}_{W^{2,\infty}}$, $\norm{W}_{W^{1,\infty}}$ and $\norm{F}_\infty$. Neither depends on $\gamma$, $\mu$, $\rho$, $\hO$, the pressure, or $\eps$.
\end{theorem}

\begin{proof}
Take $v$ from Lemma~\ref{lb:assemble} in Corollary~\ref{lb:abstract}, and use Lemma~\ref{lb:expand}:
\[
\Lambda_0\norm{\nabla e_u}\ \ge\ \frac{\ell_W(v)}{\norm{\nabla v}}-C\hG^{5/2}-C\hG^2 .
\]
Lemma~\ref{lb:assemble} bounds the first term from below.
\end{proof}

\begin{corollary}[Sharpness of Theorem~\ref{thm:D} and of Theorem~\ref{thm:A} at $G=0$]\label{lb:cor}
Assume (M0)--(M2) (so (M3) holds for small $\hG$ by Theorem~\ref{m3:thm}), $W\not\equiv0$, $\gamma\in[\gamma_2,\gamma_{\max}]$ and $\eps\le C\hG^{3/2}$ (for instance $\hO^k\le\hG^{3/2}$ and the flux condition of Theorem~\ref{thm:D}). Then for $\CNS$, and for $\GS$ when $p$ is constant on $\Gamma$,
\[
c\,\hG^{3/2}\ \le\ \norm{\nabla(\ut-u_h)}\ \le\ \tnorm{\ut-u_h}\ \le\ C\,\hG^{3/2}.
\]
So the $H^1$ rate is exactly $3/2$. For $\GS$ with $G>0$, Theorems~\ref{thm:B} and~\ref{lb:thm} together give
\[
\norm{\nabla(\ut-u_h)}\ \ge\ c\max\bigl((\hG/\gamma)G,\ \hG^{3/2}\bigr)
\]
in the regime of Theorem~\ref{thm:B}.
\end{corollary}

\begin{remark}
\begin{enumerate}
\item \emph{Only the chordal bump matters.} The lower bound is carried by the part of $\ut|_{\Gh}$ that oscillates on each chord, $-\frac12s^2W_e$ minus its mean. It has amplitude $\hG^2$ and wavelength $\hG$, hence $H^{1/2}(\Gh)$-size $\hG^{3/2}$. No discrete velocity that is Nitsche-consistent on $Y_h$ can reproduce it. The mean part $\frac{\ell^2}{12}W_e$ is smooth along $\Gh$ and contributes only $O(\hG^2)$.
\item \emph{The terms that do not carry it.} The transposed-gradient term and the penalty term of $\Gc$ vanish identically on $Y_h$, as does every pressure term. The lower bound therefore holds for $\CNS$ even though its pressure error is itself $O(\hG^{3/2})$.
\item \emph{A method-independent statement.} Corollary~\ref{lb:abstract} holds for \emph{any} discrete velocity with $a_h(u_h,v)-\ip{u_h}{\dn v}=(\ft,v)$ for all $v\in Y_h$. Every Nitsche variant with this volume form and zero boundary data, tested on $\Zh$, has this property.
\end{enumerate}
\end{remark}

\subsubsection{The route through the transposed-gradient term}\label{lb:sec:route}
The natural first attempt tests with a field whose trace is odd on each chord, and uses $\Gc(v)\approx\ip{s\,a_e}{v}$. We record what that route gives.
\begin{enumerate}
\item \emph{The relevant trace is normal, not tangential.} Since $a_e=n(m_e^\ast)\,W_e$, we have $\ip{s\,a_e}{v}=\sum_eW_e\int_es\,v\cdot n_e+O(\hG^2)\norm{v}_{L^1}$. An odd \emph{tangential} trace pairs to zero at leading order.
\item \emph{Incompressibility allows an odd normal trace.} For $v=\curl\phi$, $v\cdot n_e=-\dt\phi$, so $\int_ev\cdot n_e=0$ just says that $\phi$ takes equal values at the two ends of $e$. The identity $\dn v\cdot n_e=-\dt(v\cdot\tau_e)$ constrains the normal \emph{derivative}, not the trace. So there is no obstruction there.
\item \emph{The penalty term reinforces, at relative order $\gamma\hG$.} On $e$,
\[
\ut\cdot n_e=d\,W(x^\ast)\,\tau(x^\ast)\cdot n_e+O(\hG^4)
\qquad\text{and}\qquad
(\nabla\ut)^{\mathsf T}n_e\cdot n_e=W(x^\ast)\,\tau(x^\ast)\cdot n_e+O(\hG^2),
\]
where $\tau(x^\ast)\cdot n_e=\pm s+O(s^3)$. So on the normal trace the penalty term multiplies the transposed term pointwise by $1+(\mu/h)d(s)$, with $0\le(\mu/h)d\le\gamma\hG/8$. It is lower order for fixed $\gamma$, of the same sign, and becomes leading once $\gamma\gtrsim\hG^{-1}$. On the tangential trace the penalty pairs $\frac12(\frac{\ell^2}4-s^2)W_e$ with $v\cdot\tau_e$. This is not lower order unless $v\cdot\tau_e$ is odd or zero; choosing $\dn\phi=0$ on $\Gh$ removes it.
\item \emph{The Nitsche symmetric term is \emph{not} lower order.} $-\ip{\ut}{\dn v}$ pairs the chordal bump with $\dn v\cdot\tau_e$, which for a first-layer field of amplitude $A$ is of size $A/\hG$. It has the same order $\hG A$ as the transposed term. This is the term Theorem~\ref{lb:thm} uses.
\item \emph{Where the route fails: the pressure.} A test field with $v\cdot n\neq0$ lies outside $Y_h$. For $\CNS$ the error equation then contains $\ip{\xi-\bar\xi}{v\cdot n}$, which by Theorem~\ref{thm:D} (steps 3--4) is only bounded by
\[
C(c_0\gamma)^{-1/2}\beta^{-1}\bigl(\tnorm{e_u}+\hG^{3/2}\bigr)\tnorm v .
\]
That is the same order as the term one is trying to bound below, so it can cancel the lower bound unless $\gamma$ is large compared with constants we do not control.
\item \emph{What the route gives when it works.} For $\GS$ with $G=0$, where $|\Rc(v)|\le C\hG^2\norm{v}_{L^1(\Gh)}$ is lower order, the route gives only $\tnorm{\ut-u_h}\gtrsim\hG^{3/2}/(1+\gamma^{1/2})$ in the \emph{energy} norm. That lower bound may be carried entirely by the slip part $(\mu/h)^{1/2}\norm{e_u}_{\Gh}$. Numerically this part dominates: for $\GS$ on test A at $N=64$, $\mu=100$, it is $0.0795$ against $\norm{\nabla e_u}=0.0182$. So it says nothing about $H^1$.
\end{enumerate}

\subsubsection{Numerical checks}\label{lb:sec:num}
All computations use $k=4$ and the meshes of Section~\ref{sec:numerics}. The script is \texttt{code/}\allowbreak\texttt{lower\_bound\_tests.py}.

\begin{table}[h]
\centering
\caption{(M3) on the production meshes: star constants $\kappa_z:=\mathcal M(v_z)/|e_z|^2$, with $\norm{\nabla v_z}=1$ and $v_z$ supported in the three triangles at $z$ (floating-point star computation; the exact witness of Section~\ref{m3:sec} reproduces every entry). All production levels $N=8,\dots,256$ were checked; selected rows are shown.}\label{lb:tab:stars}
\begin{tabular}{rcccc}
\toprule
$N$ & $\dim Y^1(\omega_z)$ & $\min_z\kappa_z$ & mean & $\max_z\kappa_z$\\
\midrule
8 & 1 & 0.00288 & 0.0055 & 0.0080\\
16 & 1 & 0.00555 & 0.0109 & 0.0179\\
64 & 1 & 0.00366 & 0.0145 & 0.0252\\
256 & 1 & 0.00303 & 0.0149 & 0.0260\\
1024 & 1 & 0.00288 & 0.0149 & 0.0261\\
\bottomrule
\end{tabular}
\end{table}
\cert{code/lower\_bound\_tests.py stars 8,12,\dots,1024}{logs/lower\_bound\_stars.log}
The minimum is attained in the diagonal directions of the square, where the first-layer triangles are most elongated.

\begin{table}[h]
\centering
\caption{The assembled field $v=\sum_zW(z)v_z$ of Lemma~\ref{lb:assemble}: $\ip{\ut}{\dn v}/(\norm{\nabla v}\hG^{3/2})$, with the leading term $\ell_W(v)$ in parentheses.}\label{lb:tab:field}
\begin{tabular}{rcccc}
\toprule
$N$ & test A & rate & test B$_1$ & rate\\
\midrule
32 & 0.0606 (0.0605) & -- & 0.0184 (0.0181) & --\\
64 & 0.0651 (0.0651) & 1.39 & 0.0204 (0.0203) & 1.35\\
128 & 0.0665 (0.0665) & 1.47 & 0.0209 (0.0209) & 1.46\\
256 & 0.0669 (0.0669) & 1.49 & 0.0211 (0.0211) & 1.49\\
512 & 0.0670 (0.0670) & 1.50 & 0.0212 (0.0212) & 1.50\\
\bottomrule
\end{tabular}
\end{table}
\cert{code/lower\_bound\_tests.py field 16,\dots,512}{logs/lower\_bound\_field.log}

On the full spaces, for test A with $\mu=100$ and $N=16\to64$ (\texttt{logs/lower\_bound\_global.log}):
\begin{itemize}
\item The exact dual norm of $v\mapsto\ip{\ut}{\dn v}$ on $Y_h^1$ has rates $1.12,1.27,1.35,1.40$; its ratio to $\hG^{3/2}$ increases from $0.086$ to $0.115$.
\item It is between $14\%$ and $24\%$ of the actual $\CNS$ error $\norm{\nabla(\ut-u_h)}$.
\item On $\Zh$, the energy-dual norms of $\Gc$ and of its transposed part decay at rates $1.70\to1.54$ and $1.61\to1.51$. So Lemma~\ref{lem:transposed} is sharp on $\Zh$ as well as on $\VR$.
\item The $\norm{\nabla\cdot}$-dual norm of $\Gc$ on $\Zh$ decays only like $\hG^{1}$. It is dominated by the penalty part, which is of size $\gamma\hG$ in that norm. The $\hG^{3/2}$ upper bound genuinely needs the mesh-dependent norm.
\end{itemize}
